Berechnen Sie die Wegintegrale der folgenden Funktionen entlang der
angegebenen Pfade, die sie Sie in der vorangegangenen Aufgabe gezeichnet
haben.
\begin{teilaufgaben}
\item Für $f(z) = z^3$ berechnen Sie $\int_{\gamma_4}f(z)\,dz$.
\item Für $f(z) = \bar{z}$ berechnen Sie $\int_{\gamma_2}f(z)\,dz$ und
$\int_{\gamma_3}f(z)\,dz$.
Warum sind die beiden Integrale verschieden?
\item Für $f(z) = 1/z$ berechnen Sie $\int_{\gamma_2}f(z)\,dz$ und
$\int_{\gamma_3}f(z)\,dz$.
Warum sind die beiden Integrale verschieden?
\end{teilaufgaben}

\begin{loesung}
\begin{teilaufgaben}
\item Für die Funktion $f(z) = z^3$ sind die Wegintegrale
\begin{align*}
%\int_{\gamma_1} f(z)\,dz
%&=
%\int_0^1 (1+it)^3\, i\,dt
%\\
%&=
%\int_0^1 i - 3t - 3it^2 + t^3\,dt
%=
%\biggl[
%it - \frac{3}{2}t^2 - it^3 +\frac{1}{4}t^4
%\biggr]_0^1
%\\
%&=
%i - \frac{3}{2} -i + \frac{1}{4}
%=
%\frac{7}{4}.
%\\
%%
%% Integral über gamma_2
%%
%\int_{\gamma_2} f(z)\,dz
%&=
%\int_0^1 (e^{-\pi i t})^3 (-\pi i) e^{-\pi it}\,dt
%=
%-\pi i
%\int_0^1 e^{-4\pi i t}\,dt
%\\
%&=
%-\pi i
%\Bigl[
%\frac{1}{-4\pi i} e^{-4\pi it}
%\Bigr]_0^1
%=
%\frac{1}{4}
%\Bigl(e^{-4\pi i} - 1\Bigr)
%=
%\frac{1}{4}(1-1)
%=
%0.
%\\
%%
%% Integral über gamma_3
%%
%\int_{\gamma_3} f(z)\,dz
%&=
%\int_0^1 (e^{\pi i t})^3 (\pi i) e^{\pi it}\,dt
%=
%\pi i
%\int_0^1 e^{4\pi i t}\,dt
%\\
%&=
%\pi i
%\Bigl[
%\frac{1}{4\pi i} e^{4\pi it}
%\Bigr]_0^1
%=
%\frac{1}{4}
%\Bigl(e^{4\pi i} - 1\Bigr)
%=
%\frac{1}{4}(1-1)
%=
%0.
%\\
\int_{\gamma_4} f(z)\,dz
&=
\int_0^1 
(1+it+t^2)^3 (i+2t)
\,dt
=
\int_0^1 
2t^3+it^2+2it+2t+i-1
\,dt
\\
&=
\biggl[
\frac{t^4}{2}
+
\frac{it^3}{3}
+
it^2
+
t^2
+it
-
t
\biggr]_0^1
=
\frac{14i+3}{6}.
\end{align*}
\item Für die Funktion $f(z) = \bar{z}$ sind die Wegintegrale
\begin{align*}
%\int_{\gamma_1} f(z)\,dz
%&=
%\int_0^1
%\overline{(1+it)}
%i
%\,dt
%=
%\int_0^1
%i+t
%\,dt
%=
%\biggl[ it+\frac{t^2}{2} \biggr]_0^1
%=
%i+\frac12
%\\
\int_{\gamma_2} f(z)\,dz
&=
\int_0^1
\overline{e^{-\pi it}}
(-\pi i)e^{-\pi it}
\,dt
=
\int_0^1
e^{\pi it}
(-\pi i)e^{-\pi it}
\,dt
=
-\pi i
\int_0^1
\,dt
=
-\pi i.
\\
\int_{\gamma_3} f(z)\,dz
&=
\int_0^1
\overline{e^{\pi it}}
\pi ie^{\pi it}
\,dt
=
\int_0^1
e^{-\pi it}
\pi ie^{\pi it}
\,dt
=
\pi i
\int_0^1
\,dt
=
\pi i.
%\\
%\int_{\gamma_4} f(z)\,dz
%&=
%\int_0^1
%\overline{(1+it+t^2)}
%(i+2t)
%\,dt
%=
%\int_0^1
%(1-it+t^2)
%(i+2t)
%\,dt
%\\
%&=
%\int_0^1
%(1-it+t^2)
%(i+2t)
%\,dt
%=
%\int_0^1
%2t^3-it^2+3t+i
%\,dt.
%\\
%&=
%\biggl[
%\frac{t^4}{2}-\frac{it^3}{3}+\frac{3t^2}{2}+it
%\biggr]_0^1
%=
%\frac{2i+6}{3}.
\end{align*}
Die beiden Integrale sind verschieden, weil die Funktion $z\mapsto\bar{z}$
nicht holomorph ist.
\item %Für die Funktion $f(z) = 1/z$ sind die verlangten Wegintegrale
%\begin{align*}
%\int_{\gamma_1} f(z)\,dz
%&=
%\int_0^1 \frac{1}{1+it}\,i\,dt
%=
%\int_0^1 \frac{1-it}{1+t^2}\,i\,dt
%=
%\int_0^1 \frac{i+t}{1+t^2}\,dt
%\\
%&=
%\biggl[
%\frac12 \log(1+t^2) + i\arctan(t)
%\biggr]_0^1
%=
%\frac12\log 2 + \frac{i\pi}4,
%\intertext{
Wegen $1/e^{i\pi t}=e^{-i\pi t}=\overline{e^{i\pi t}}$ stimmt
die Funktion $z\mapsto 1/z$ auf dem Einheitskreis mit $z\mapsto\bar{z}$ 
überein, ergibt also die gleichen Integrale wie in Teilaufgabe b).
%Das verbleibende Integral über $\gamma_4$ ist
%}
%\int_{\gamma_4} f(z)\,dz
%&=
%\int_0^1 \frac{1}{1+it+t^2} (i+2t)\,dt
%=
%\biggl[
%\log(t^2+it+1)
%\biggr]_0^1
%\\
%&=
%\log(1+i) - \log(1)
%=
%\log(1+i)
%=
%\frac12
%\log(2)
%+
%\frac{i\pi}{4}.
Die beiden Integrale sind verschieden, weil die Funktion $z\mapsto1/z$
ist zwar holomoroph, aber nicht im Punkt $z=0$, wo sie einen Pol hat.
\qedhere
%\end{align*}
\end{teilaufgaben}
\end{loesung}
